% ===========================================================================
\chapter{Long conversations}\label{ch:long}

Every request sends the whole conversation. On a long task, with dozens of
files read and commands run, the conversation grows, and with it three
costs. Each request takes longer. Each costs more: the tenth request reads
everything the first nine did, so the total read grows with the
\emph{square} of the number of replies, and caching (Chapter~\ref{ch:claude})
only softens that. And a model can only read so much, its \emph{context
window}; a conversation that outgrows it can't be sent at all. Long before
that point, a model working through a huge conversation also tends to work
less well: the one detail that matters is buried among a hundred that
don't.

So an agent needs a plan for when the conversation gets long.

\section{What not to do}

The tempting fix is to edit the conversation: drop old tool results (who
needs a file listing from forty replies ago?), or shorten them. Don't. We
decided in Chapter~\ref{ch:conversation} that the model's replies go back
exactly as they came, and the same goes for the rest of the conversation
once it has been sent. A model that thinks signs its thinking, and the
API checks that the conversation the thinking belongs to has not been
changed since; an edited history can be rejected outright. And the prompt
cache is keyed on the exact text of the conversation, so editing an early
message makes every later one a cache miss. The conversation should only
ever grow at the end.

What we already do helps: every tool result is cut to about 5,000 tokens
(Chapter~\ref{ch:tools}), and the descriptions steer the model towards
searching rather than reading everything.

\section{Summing up and starting afresh}

When the conversation passes a budget, the agent asks the model to sum up:
the task, what it has found, what it has changed, what is left. Then it
starts a new conversation, whose first message is the task and that
summary, and carries on. This is what a programmer does at the end of a
long day: write a note for tomorrow, and start fresh from the note.

The request for the summary, and the start of the new conversation, go in
\code{context.nex}:

\begin{append}{context.nex}
-- What the agent asks for when the conversation has grown too long.
function summary_request(): String do
  result := "This conversation is getting long, so it will be cut short here. "
    + "Don't call any tools now. Write a summary for yourself to carry on from: "
    + "what the task is, what you have found out (the files and the lines that "
    + "matter), what you have changed so far, and what is left to do."
end

-- A new conversation that starts from `summary`, for carrying on with `task`.
function fresh_start(task, summary: String): Conversation
require
  task_given: task /= ""
do
  result := create Conversation.make()
  result.add_user_text("You are part way through a task. The conversation so far "
    + "was too long to keep, so here is the task and your own summary of the work "
    + "so far. Carry on from there.\n\nThe task:\n" + task
    + "\n\nYour summary:\n" + (when summary = "" then "(there is no summary)" else summary end))
ensure
  models_turn: result.models_turn()
end
\end{append}

The summary is written by the model, for the model: it knows best what it
will need. The summary request names what to keep, so that nothing
essential is lost, such as which files were already changed.

\section{In the agent}

The agent gets a budget, in tokens. 200,000 is a sensible default: far
below the context window of current models, but already a long and
expensive conversation. Add the field, its default, and a way to change it:

\begin{tightaddbefore}{agent.nex}{after max_replies}{|  set_max_replies(n: Integer)}
  context_limit: Integer       -- the longest conversation, in tokens
\end{tightaddbefore}

\begin{tightaddbefore}{agent.nex}{at the end of make}{  end||feature|  model: Model}
    context_limit := 200000
\end{tightaddbefore}

\begin{addbefore}{agent.nex}{before run}{  -- Work on `task` until the model says it is done, or the agent stops}
  set_context_limit(tokens: Integer)
  require
    positive: tokens > 0
  do
    context_limit := tokens
  end
\end{addbefore}

How long is the conversation? The API tells us with every reply: its
\code{context_tokens} (Chapter~\ref{ch:conversation}) is how much the request
held. The right moment to check is just after the tool results have been
added, when the conversation is complete and it is the model's turn. In
\code{run}, right after \code{conversation.add_tool_results(results)}, add:

\begin{tightaddbefore}{agent.nex}{in run, after the tool results}{        else|          ending := ending_for(reply.stop_reason())}
          if reply.context_tokens() > context_limit and model.can_reply() then
            make_room(task, reply.context_tokens())
          end
\end{tightaddbefore}

\code{make_room} adds the summary request to the conversation (it joins the
message with the tool results, as \code{add_user_text} promised in
Chapter~\ref{ch:conversation}), gets the summary, and replaces the
conversation with a fresh one. The summary costs one more reply, which is
counted in the usage. The model was told not to call tools; if it does
anyway, the calls are simply dropped with the old conversation.

\begin{addbefore}{agent.nex}{in the private features, before ending_for}{    -- What a reply that asks for no tools means for the task.}
    -- The conversation is too long: have the model sum it up, and carry
    -- on in a new conversation that starts from the summary.
    make_room(task: String, tokens: Integer)
    require
      models_turn: conversation.models_turn()
    do
      reporter.on_making_room(tokens)
      conversation.add_user_text(summary_request())
      let reply: Reply := model.reply(create Request.make(system, conversation, definitions()))
      usage.add(reply)
      conversation := fresh_start(task, reply.text())
    ensure
      models_turn: conversation.models_turn()
    end
\end{addbefore}

And the reporter says when it happens:

\begin{addbefore}{agent.nex}{in Reporter, before on_finish}{  on_finish(outcome: Run_Result, usage: Usage) do}
  on_making_room(tokens: Integer) do
    say("  (the conversation is " + grouped(tokens) + " tokens long: summing up and starting afresh)")
  end
\end{addbefore}

To see it work without a long conversation, we give a script replies that
claim to have read a lot, and set a small budget:

\begin{tryit}{try_long.nex}
intern Agent

let ws: Workspace := fresh_workspace(user_path("sample_project"), user_path("scratch"))
let model: Scripted_Model := create Scripted_Model.make([
  create Reply.from_json(json_value('{"content": [{"type": "tool_use", "id": "t1", "name": "list_files", "input": {}}], "stop_reason": "tool_use", "usage": {"input_tokens": 3000, "output_tokens": 20}}')),
  say("The task: explain the project. Found: four files; fines are in fines.nex. Changed: nothing. Left: explain."),
  say("It is a lending library: loans, and fines for late books.")
])
let agent: Agent := create Agent.make(model, standard_tools(ws), create Allow_Everything.make(),
  create Reporter.make(false, false))
agent.set_context_limit(1000)
agent.run("Explain this project.")
console.print_line("\nThe new conversation starts:")
console.print_line(model.sent.get(2))
\end{tryit}
\tryoutput{try_long}

After the summary, the conversation is one message long, and the work
carries on from it.

\begin{aside}
\textbf{The API can do this too.} Claude's API has a beta feature called
\emph{compaction}, which sums up a long conversation on the server, in
a form the model reads back well, and keeps it valid for thinking and
caching. It is a good choice for a production agent. Doing it ourselves,
as here, shows what is involved, and works with any model.
\end{aside}
